% Options for packages loaded elsewhere
% Options for packages loaded elsewhere
\PassOptionsToPackage{unicode}{hyperref}
\PassOptionsToPackage{hyphens}{url}
%
\documentclass[
  10pt,
  ignorenonframetext,
  aspectratio=169,
  brazilian,
]{beamer}
\newif\ifbibliography
\usepackage{pgfpages}
\setbeamertemplate{caption}[numbered]
\setbeamertemplate{caption label separator}{: }
\setbeamercolor{caption name}{fg=normal text.fg}
\beamertemplatenavigationsymbolsempty
% remove section numbering
\setbeamertemplate{part page}{
  \centering
  \begin{beamercolorbox}[sep=16pt,center]{part title}
    \usebeamerfont{part title}\insertpart\par
  \end{beamercolorbox}
}
\setbeamertemplate{section page}{
  \centering
  \begin{beamercolorbox}[sep=12pt,center]{section title}
    \usebeamerfont{section title}\insertsection\par
  \end{beamercolorbox}
}
\setbeamertemplate{subsection page}{
  \centering
  \begin{beamercolorbox}[sep=8pt,center]{subsection title}
    \usebeamerfont{subsection title}\insertsubsection\par
  \end{beamercolorbox}
}
% Prevent slide breaks in the middle of a paragraph
\widowpenalties 1 10000
\raggedbottom
\AtBeginPart{
  \frame{\partpage}
}
\AtBeginSection{
  \ifbibliography
  \else
    \frame{\sectionpage}
  \fi
}
\AtBeginSubsection{
  \frame{\subsectionpage}
}
\usepackage{iftex}
\ifPDFTeX
  \usepackage[T1]{fontenc}
  \usepackage[utf8]{inputenc}
  \usepackage{textcomp} % provide euro and other symbols
\else % if luatex or xetex
  \usepackage{unicode-math} % this also loads fontspec
  \defaultfontfeatures{Scale=MatchLowercase}
  \defaultfontfeatures[\rmfamily]{Ligatures=TeX,Scale=1}
\fi
\usepackage{lmodern}

\ifPDFTeX\else
  % xetex/luatex font selection
\fi
% Use upquote if available, for straight quotes in verbatim environments
\IfFileExists{upquote.sty}{\usepackage{upquote}}{}
\IfFileExists{microtype.sty}{% use microtype if available
  \usepackage[]{microtype}
  \UseMicrotypeSet[protrusion]{basicmath} % disable protrusion for tt fonts
}{}
\makeatletter
\@ifundefined{KOMAClassName}{% if non-KOMA class
  \IfFileExists{parskip.sty}{%
    \usepackage{parskip}
  }{% else
    \setlength{\parindent}{0pt}
    \setlength{\parskip}{6pt plus 2pt minus 1pt}}
}{% if KOMA class
  \KOMAoptions{parskip=half}}
\makeatother


\usepackage{longtable,booktabs,array}
\newcounter{none} % for unnumbered tables
\usepackage{calc} % for calculating minipage widths
\usepackage{caption}
% Make caption package work with longtable
\makeatletter
\def\fnum@table{\tablename~\thetable}
\makeatother
\usepackage{graphicx}
\makeatletter
\newsavebox\pandoc@box
\newcommand*\pandocbounded[1]{% scales image to fit in text height/width
  \sbox\pandoc@box{#1}%
  \Gscale@div\@tempa{\textheight}{\dimexpr\ht\pandoc@box+\dp\pandoc@box\relax}%
  \Gscale@div\@tempb{\linewidth}{\wd\pandoc@box}%
  \ifdim\@tempb\p@<\@tempa\p@\let\@tempa\@tempb\fi% select the smaller of both
  \ifdim\@tempa\p@<\p@\scalebox{\@tempa}{\usebox\pandoc@box}%
  \else\usebox{\pandoc@box}%
  \fi%
}
% Set default figure placement to htbp
\def\fps@figure{htbp}
\makeatother



\ifLuaTeX
\usepackage[bidi=basic,shorthands=off]{babel}
\else
\usepackage[bidi=default,shorthands=off]{babel}
\fi
\ifLuaTeX
  \usepackage{selnolig} % disable illegal ligatures
\fi


\setlength{\emergencystretch}{3em} % prevent overfull lines

\providecommand{\tightlist}{%
  \setlength{\itemsep}{0pt}\setlength{\parskip}{0pt}}



 


\makeatletter
\@ifpackageloaded{caption}{}{\usepackage{caption}}
\AtBeginDocument{%
\ifdefined\contentsname
  \renewcommand*\contentsname{Índice}
\else
  \newcommand\contentsname{Índice}
\fi
\ifdefined\listfigurename
  \renewcommand*\listfigurename{Lista de Figuras}
\else
  \newcommand\listfigurename{Lista de Figuras}
\fi
\ifdefined\listtablename
  \renewcommand*\listtablename{Lista de Tabelas}
\else
  \newcommand\listtablename{Lista de Tabelas}
\fi
\ifdefined\figurename
  \renewcommand*\figurename{Figura}
\else
  \newcommand\figurename{Figura}
\fi
\ifdefined\tablename
  \renewcommand*\tablename{Tabela}
\else
  \newcommand\tablename{Tabela}
\fi
}
\@ifpackageloaded{float}{}{\usepackage{float}}
\floatstyle{ruled}
\@ifundefined{c@chapter}{\newfloat{codelisting}{h}{lop}}{\newfloat{codelisting}{h}{lop}[chapter]}
\floatname{codelisting}{Listagem}
\newcommand*\listoflistings{\listof{codelisting}{Lista de Listagens}}
\makeatother
\makeatletter
\makeatother
\makeatletter
\@ifpackageloaded{caption}{}{\usepackage{caption}}
\@ifpackageloaded{subcaption}{}{\usepackage{subcaption}}
\makeatother

\usepackage{bookmark}
\IfFileExists{xurl.sty}{\usepackage{xurl}}{} % add URL line breaks if available
\urlstyle{same}
% fallback for those not using the hyperref driver hyperxmp:
\makeatletter
\@ifundefined{xmpquote}{\newcommand{\xmpquote}[1]{#1}}{}
\makeatother
\hypersetup{
  pdftitle={Parte I : Aprendizado Supervisionado},
  pdflang={pt-BR},
  hidelinks,
  pdfcreator={LaTeX via pandoc}}


\title{\texorpdfstring{Parte I : Aprendizado
Supervisionado}{Parte I : Aprendizado Supervisionado}}
\subtitle{\texorpdfstring{Alta Dimensão, Decisão, Avaliação e
Interpretabilidade}{Alta Dimensão, Decisão, Avaliação e Interpretabilidade}}
\author{}
\date{}

\begin{document}
\frame{\titlepage}


\begin{frame}
\end{frame}

\begin{frame}{Alta dimensionalidade: qual é o problema?}
\protect\phantomsection\label{alta-dimensionalidade-qual-uxe9-o-problema}
Queremos aprender a relação

\[
Y=f(\mathbf X)+\varepsilon,\qquad \mathbf X\in\mathbb R^p,
\quad \mathbb E(\varepsilon\mid\mathbf X)=0,
\]

utilizando \(n\) observações e \(p\) variáveis de entrada.

Com muitas entradas, precisamos de dados suficientes para distinguir o
sinal do ruído. A dificuldade pode aparecer mesmo quando \(p<n\).

\textbf{Percurso:} estimar em alta dimensão, explorar estrutura e
regularização, decidir com custos, avaliar e interpretar o modelo.
\end{frame}

\begin{frame}{A maldição da dimensionalidade}
\protect\phantomsection\label{a-maldiuxe7uxe3o-da-dimensionalidade}
Imagine observações uniformes em \([0,1]^p\). Uma região com comprimento
\(h\) em cada coordenada ocupa a fração \(h^p\) do volume.

Para manter \(k\) observações \textbf{em média} nessa região:

\[
\boxed{nh^p=k\quad\Longrightarrow\quad n=\frac{k}{h^p}.}
\]

{\def\LTcaptype{none} % do not increment counter
\begin{longtable}[]{@{}lrrr@{}}
\toprule\noalign{}
Com \(h=0{,}2\) e \(k=20\) & \(p=1\) & \(p=3\) & \(p=5\) \\
\midrule\noalign{}
\endhead
Amostra necessária, em média & 100 & 2.500 & 62.500 \\
\bottomrule\noalign{}
\end{longtable}
}

Manter vizinhanças pequenas exige muito mais dados. Isso afeta
especialmente KNN e regressão local.

\textbf{Esta conta ilustra cobertura do espaço sob uniformidade. Não é
uma regra geral de tamanho amostral.}
\end{frame}

\begin{frame}{O que é uma taxa de convergência?}
\protect\phantomsection\label{o-que-uxe9-uma-taxa-de-converguxeancia}
O risco de estimação pode ser medido por

\[
R_n=\mathbb E[(\hat f_n(\mathbf X)-f(\mathbf X))^2],
\]

com esperança sobre o treinamento e uma nova entrada. Ele não inclui o
ruído irredutível da resposta.

\textbf{Como ilustração}, se \(R_n\approx Cn^{-a}\), então
\(R_{4n}/R_n\approx4^{-a}\). \(C\) é uma constante de escala, \(a\) é o
expoente e 4 é o fator de aumento da amostra.

{\def\LTcaptype{none} % do not increment counter
\begin{longtable}[]{@{}rrr@{}}
\toprule\noalign{}
Expoente & Erro relativo & Redução ilustrativa \\
\midrule\noalign{}
\endhead
\(a=1\) & \(0{,}25\) & 75\% \\
\(a=1/2\) & \(0{,}50\) & 50\% \\
\(a=1/6\) & \(0{,}79\) & 21\% \\
\bottomrule\noalign{}
\end{longtable}
}

O modelo de potência não é universal. Uma relação de ordem
\(R_n\asymp n^{-a}\), sozinha, não garante esses percentuais.
\end{frame}

\begin{frame}{Taxas de convergência e dimensionalidade}
\protect\phantomsection\label{taxas-de-converguxeancia-e-dimensionalidade}
Para classes gerais de funções suaves em \(p\) dimensões, uma taxa
minimax típica do erro quadrático integrado é

\[
\boxed{R_n^*\asymp n^{-\frac{2\alpha}{2\alpha+p}}.}
\]

\begin{itemize}
\tightlist
\item
  \(\alpha\): suavidade da função. Valores maiores impõem maior
  regularidade.
\item
  \(R_n^*\): melhor risco possível no pior caso da classe considerada.
\item
  \(\asymp\): mesma ordem, desconsiderando constantes multiplicativas.
\end{itemize}

\textbf{Aumentar \(p\) reduz o expoente e torna a aprendizagem mais
lenta.}

A taxa exige condições sobre a classe de funções, o desenho e o ruído.
Mede estimação de \(f\), sem incluir o ruído irredutível de uma nova
resposta.

\note{Formalmente, R\_n\^{}*=inf\_\{estimadores\} sup\_\{f em F\_alpha\}
E{[}(f\_estimado(X)-f(X))²{]}, com esperança sobre treinamento e nova
entrada. Taxa minimax usual sob condições de regularidade, dimensão fixa
e classe suave. Fonte: Wasserman, Nonparametric Regression.}
\end{frame}

\begin{frame}{Quanto ganhamos aumentando a amostra?}
\protect\phantomsection\label{quanto-ganhamos-aumentando-a-amostra}
Fixando suavidade \(\alpha=2\), a taxa é \(n^{-4/(4+p)}\).

{\def\LTcaptype{none} % do not increment counter
\begin{longtable}[]{@{}rrr@{}}
\toprule\noalign{}
Dimensão \(p\) & Taxa & Fator do termo principal com \(4n\) \\
\midrule\noalign{}
\endhead
1 & \(n^{-4/5}\) & \(0{,}33\) \\
4 & \(n^{-1/2}\) & \(0{,}50\) \\
10 & \(n^{-2/7}\) & \(0{,}67\) \\
20 & \(n^{-1/6}\) & \(0{,}79\) \\
\bottomrule\noalign{}
\end{longtable}
}

Em uma dimensão, o termo cai para cerca de 33\% do valor anterior. Em 20
dimensões, permanece em cerca de 79\%.

\textbf{Para reduzir esse termo à metade em \(p=20\), seria necessário
multiplicar \(n\) por \(2^6=64\).} São comparações da taxa dentro de
cada cenário, sem considerar diferenças nas constantes.
\end{frame}

\begin{frame}{A redução do erro fica mais lenta}
\protect\phantomsection\label{a-reduuxe7uxe3o-do-erro-fica-mais-lenta}
\pandocbounded{\includegraphics[keepaspectratio]{aula_altas_dimensoes_est335_revisada_V2_files/figure-beamer/unnamed-chunk-1-1.png}}

Todas as curvas começam em 1 para comparar \textbf{velocidades de
redução}, não níveis absolutos de erro. São curvas teóricas com
\(\alpha=2\), sem resultados de ajuste a dados.
\end{frame}

\begin{frame}{Como tornar o problema mais simples?}
\protect\phantomsection\label{como-tornar-o-problema-mais-simples}
A taxa anterior considera uma classe ampla de funções em \(p\)
dimensões.

Podemos explorar hipóteses que reduzam a complexidade relevante:

{\def\LTcaptype{none} % do not increment counter
\begin{longtable}[]{@{}
  >{\raggedright\arraybackslash}p{(\linewidth - 4\tabcolsep) * \real{0.3333}}
  >{\raggedright\arraybackslash}p{(\linewidth - 4\tabcolsep) * \real{0.3333}}
  >{\raggedright\arraybackslash}p{(\linewidth - 4\tabcolsep) * \real{0.3333}}@{}}
\toprule\noalign{}
\begin{minipage}[b]{\linewidth}\raggedright
Hipótese
\end{minipage} & \begin{minipage}[b]{\linewidth}\raggedright
Ideia
\end{minipage} & \begin{minipage}[b]{\linewidth}\raggedright
O que procuramos
\end{minipage} \\
\midrule\noalign{}
\endhead
Esparsidade & Poucas entradas são necessárias à função de regressão &
Selecionar efeitos relevantes \\
Redundância & Muitas entradas carregam informação compartilhada &
Combinar ou representar informação comum \\
\bottomrule\noalign{}
\end{longtable}
}

\textbf{Ter muitas variáveis observadas não implica que a resposta
dependa de todas de forma independente e complexa.}
\end{frame}

\begin{frame}{Esparsidade: poucos efeitos ativos}
\protect\phantomsection\label{esparsidade-poucos-efeitos-ativos}
Suponha \(p=1.000\), mas

\[
Y=3X_2-2X_{17}+1{,}5X_{204}+\varepsilon.
\]

A função de regressão utiliza apenas \(s=3\) entradas: \(s\ll p\).

Se soubéssemos quais são, o ajuste seria muito mais simples. Como não
sabemos, precisamos encontrá-las entre as 1.000 candidatas.

\textbf{Lasso pode explorar esparsidade dos coeficientes.} A hipótese
também pode ocorrer em funções não lineares, mas lasso nas entradas
originais representa apenas efeitos lineares.

Esparsidade do modelo significa poucos efeitos ativos, sem exigir muitos
zeros na matriz de dados.
\end{frame}

\begin{frame}{Redundância: muitas medições do mesmo fenômeno}
\protect\phantomsection\label{redunduxe2ncia-muitas-mediuxe7uxf5es-do-mesmo-fenuxf4meno}
Imagine 100 sensores que medem uma condição térmica comum:

\[
X_j=a_jZ+\eta_j,\qquad j=1,\ldots,100.
\]

\(Z\) é a condição compartilhada, \(a_j\) é a sensibilidade do sensor e
\(\eta_j\) é o ruído específico da medição.

As entradas podem ser fortemente correlacionadas. Se \(Y\) depende
principalmente de \(Z\), a informação preditiva está associada a uma
estrutura simples.

\textbf{Ridge pode estabilizar o uso conjunto dessas medições.} Uma RNA
pode aprender combinações, mas precisa de dados e controle de
complexidade.

Com ruído em todos os sensores, a estrutura de baixa dimensão é
aproximada. Remover medições também pode eliminar informação útil.
\end{frame}

\begin{frame}{Esparsidade e redundância são diferentes}
\protect\phantomsection\label{esparsidade-e-redunduxe2ncia-suxe3o-diferentes}
{\def\LTcaptype{none} % do not increment counter
\begin{longtable}[]{@{}
  >{\raggedright\arraybackslash}p{(\linewidth - 4\tabcolsep) * \real{0.3333}}
  >{\raggedright\arraybackslash}p{(\linewidth - 4\tabcolsep) * \real{0.3333}}
  >{\raggedright\arraybackslash}p{(\linewidth - 4\tabcolsep) * \real{0.3333}}@{}}
\toprule\noalign{}
\begin{minipage}[b]{\linewidth}\raggedright
Aspecto
\end{minipage} & \begin{minipage}[b]{\linewidth}\raggedright
Esparsidade
\end{minipage} & \begin{minipage}[b]{\linewidth}\raggedright
Redundância
\end{minipage} \\
\midrule\noalign{}
\endhead
Estrutura & Poucos efeitos necessários & Informação compartilhada \\
Exemplo & Três entre 1.000 entradas & 100 sensores de um fator \\
Estratégia & Selecionar ou contrair efeitos & Estabilizar ou combinar
medições \\
Cuidado & Achar quais efeitos importam & Preservar o sinal para \(Y\) \\
\bottomrule\noalign{}
\end{longtable}
}

As duas hipóteses podem coexistir: muitos grupos de entradas
semelhantes, mas apenas alguns grupos relacionados à resposta.

\textbf{Uma combinação preditiva de dimensão um pode usar todas as
entradas. Dimensão pequena não exige coeficientes esparsos na
representação original.}
\end{frame}

\begin{frame}{Estrutura e tamanho da amostra}
\protect\phantomsection\label{estrutura-e-tamanho-da-amostra}
Na regressão linear esparsa, sob condições adequadas, limites de erro
preditivo incluem termos da ordem

\[
\boxed{\sigma^2\frac{s\log p}{n}.}
\]

\begin{itemize}
\tightlist
\item
  \(s\): quantidade de efeitos ativos.
\item
  \(\log p\): custo de procurar efeitos entre muitos candidatos.
\item
  \(\sigma^2\): variância do ruído.
\end{itemize}

Isso ajuda a entender por que aprender com \(p>n\) pode ser viável.

\textbf{A expressão não determina uma amostra mínima universal.}
Correlação, intensidade do sinal e precisão desejada também importam.
Prever bem e identificar exatamente as entradas ativas são objetivos
diferentes.

\note{Limite típico para regressão linear esparsa com regularização
calibrada, controle de caudas e condições de compatibilidade/autovalores
restritos. Não é uma taxa para qualquer função não linear esparsa.
Fonte: Gautier e Tsybakov.}
\end{frame}

\begin{frame}{Ridge: estabilidade com muitas entradas}
\protect\phantomsection\label{ridge-estabilidade-com-muitas-entradas}
Ridge minimiza o erro de ajuste acrescido de uma penalização quadrática:

\[
\sum_{i=1}^{n}(y_i-\beta_0-\mathbf x_i^\top\boldsymbol\beta)^2
+\lambda\sum_{j=1}^{p}\beta_j^2.
\]

\textbf{Quando considerar:} entradas correlacionadas ou sinal
distribuído por muitos efeitos pequenos.

Para \(\lambda>0\), o critério produz coeficientes únicos mesmo com
\(p>n\), no ajuste com intercepto separado.

A penalização estabiliza o ajuste, mas geralmente mantém todos os
coeficientes. Uma solução única não garante boa predição: \(\lambda\)
precisa ser escolhido por validação.
\end{frame}

\begin{frame}{Lasso: regularização com seleção}
\protect\phantomsection\label{lasso-regularizauxe7uxe3o-com-seleuxe7uxe3o}
Lasso substitui a penalização quadrática pela soma dos valores
absolutos:

\[
\sum_{i=1}^{n}(y_i-\beta_0-\mathbf x_i^\top\boldsymbol\beta)^2
+\lambda\sum_{j=1}^{p}|\beta_j|.
\]

\textbf{Quando considerar:} poucos efeitos aproximadamente lineares na
representação utilizada.

O critério favorece coeficientes exatamente iguais a zero, podendo
reduzir o conjunto de entradas utilizadas.

Com entradas fortemente correlacionadas, a escolha de variáveis pode
mudar entre amostras. \textbf{Selecionar uma variável não prova que ela
seja a única portadora do sinal.}
\end{frame}

\begin{frame}{SVM linear: classificação com regularização}
\protect\phantomsection\label{svm-linear-classificauxe7uxe3o-com-regularizauxe7uxe3o}
A SVM linear utiliza uma função de decisão

\[
g(\mathbf x)=\mathbf w^\top\mathbf x+b
\]

com controle da magnitude de \(\mathbf w\) e penalização das violações
da margem.

\textbf{Quando considerar:} classificação com muitas entradas e uma
fronteira linear plausível, inclusive em problemas com \(p>n\).

O parâmetro \(C\) equilibra regularização e violações da margem. Deve
ser escolhido por validação, com padronização quando apropriada.

Um kernel não linear pode ampliar a flexibilidade, mas distâncias pouco
informativas não passam a ser úteis apenas pelo uso do kernel.
\end{frame}

\begin{frame}{RNA: possibilidade de aprender representações}
\protect\phantomsection\label{rna-possibilidade-de-aprender-representauxe7uxf5es}
Uma camada oculta calcula

\[
\mathbf a^{(1)}=\phi(W^{(1)}\mathbf x+\mathbf b^{(1)}).
\]

Ela pode combinar entradas redundantes e representar relações não
lineares.

\textbf{Quando considerar:} estrutura não linear que possa ser aprendida
com os dados disponíveis, arquitetura adequada e regularização.

\begin{itemize}
\tightlist
\item
  Com \(H\) neurônios, a primeira camada já possui \(H(p+1)\)
  parâmetros.
\item
  Weight decay penaliza pesos. Early stopping controla a duração do
  ajuste pela validação.
\end{itemize}

\textbf{Uma RNA não é automaticamente uma boa escolha para \(p\gg n\).}
Com amostra pequena, ridge e lasso são referências importantes para
comparação.
\end{frame}

\begin{frame}{Qual método considerar primeiro?}
\protect\phantomsection\label{qual-muxe9todo-considerar-primeiro}
{\def\LTcaptype{none} % do not increment counter
\begin{longtable}[]{@{}
  >{\raggedright\arraybackslash}p{(\linewidth - 2\tabcolsep) * \real{0.5000}}
  >{\raggedright\arraybackslash}p{(\linewidth - 2\tabcolsep) * \real{0.5000}}@{}}
\toprule\noalign{}
\begin{minipage}[b]{\linewidth}\raggedright
Situação
\end{minipage} & \begin{minipage}[b]{\linewidth}\raggedright
Método estudado a considerar
\end{minipage} \\
\midrule\noalign{}
\endhead
Muitas entradas correlacionadas ou muitos efeitos pequenos & Ridge \\
Poucos efeitos lineares relevantes & Lasso \\
Classificação com fronteira aproximadamente linear & SVM linear
regularizada \\
Relações não lineares e dados suficientes para aprendê-las & RNA
regularizada \\
\bottomrule\noalign{}
\end{longtable}
}

A tabela orienta comparações, sem definir um vencedor universal.

\textbf{Procedimento essencial:} estimar padronização, seleção e
hiperparâmetros dentro do treinamento/validação. Reservar o teste para a
avaliação final.
\end{frame}

\begin{frame}{Da estimação à decisão}
\protect\phantomsection\label{da-estimauxe7uxe3o-uxe0-decisuxe3o}
Em alta dimensão, regularizamos para aprender com os dados disponíveis.
Depois do ajuste, ainda precisamos responder:

{\def\LTcaptype{none} % do not increment counter
\begin{longtable}[]{@{}ll@{}}
\toprule\noalign{}
Etapa & Pergunta \\
\midrule\noalign{}
\endhead
Predizer & Qual resposta ou probabilidade o modelo estima? \\
Decidir & Que ação tomar diante dessa predição? \\
Avaliar & Como o procedimento funciona em dados novos? \\
Interpretar & Como o modelo produz suas predições? \\
\bottomrule\noalign{}
\end{longtable}
}

\textbf{Bom ajuste, boa decisão e boa explicação são propriedades
diferentes.}
\end{frame}

\begin{frame}{Probabilidade, cutoff e regra de Bayes}
\protect\phantomsection\label{probabilidade-cutoff-e-regra-de-bayes}
Na classificação binária, defina
\(\eta(\mathbf x)=P(Y=1\mid\mathbf X=\mathbf x)\).

Com custos iguais para os dois tipos de erro, a regra de Bayes escolhe

\[
g^*(\mathbf x)=I\{\eta(\mathbf x)>0{,}5\}.
\]

Na prática, substituímos \(\eta\) pela probabilidade estimada
\(\hat\eta\) e escolhemos um cutoff \(c\):

\[
\hat g_c(\mathbf x)=I\{\hat\eta(\mathbf x)>c\}.
\]

Mudar \(c\) altera as classes previstas, sem alterar as probabilidades
do modelo. Mesmo a regra ótima erra quando as classes se sobrepõem:

\[
R_{\mathrm{Bayes}}=\mathbb E[\min\{\eta(\mathbf X),1-\eta(\mathbf X)\}].
\]
\end{frame}

\begin{frame}{Custos diferentes produzem decisões diferentes}
\protect\phantomsection\label{custos-diferentes-produzem-decisuxf5es-diferentes}
Sejam \(C_{FP}\) o custo do falso positivo e \(C_{FN}\) o custo do falso
negativo. Acertos têm custo zero.

\[
R(1\mid\mathbf x)=C_{FP}[1-\eta(\mathbf x)],
\qquad R(0\mid\mathbf x)=C_{FN}\eta(\mathbf x).
\]

Escolher a ação de menor risco equivale a decidir 1 quando

\[
\boxed{\eta(\mathbf x)>c^*=\frac{C_{FP}}{C_{FP}+C_{FN}}.}
\]

\textbf{Exemplo:} \(C_{FP}=1\) e \(C_{FN}=9\) dão \(c^*=0{,}10\). Uma
probabilidade de 0,20 já justifica a decisão positiva.

A regra pressupõe probabilidades populacionais adequadas e custos
constantes. Aplicá-la a \(\hat\eta\) requer atenção à calibração,
especialmente após reamostragem ou ponderação de classes.
\end{frame}

\begin{frame}{Matriz de confusão e perguntas das métricas}
\protect\phantomsection\label{matriz-de-confusuxe3o-e-perguntas-das-muxe9tricas}
{\def\LTcaptype{none} % do not increment counter
\begin{longtable}[]{@{}lrr@{}}
\toprule\noalign{}
Classe observada & Predito 1 & Predito 0 \\
\midrule\noalign{}
\endhead
1 & Verdadeiro positivo (VP) & Falso negativo (FN) \\
0 & Falso positivo (FP) & Verdadeiro negativo (VN) \\
\bottomrule\noalign{}
\end{longtable}
}

{\def\LTcaptype{none} % do not increment counter
\begin{longtable}[]{@{}
  >{\raggedright\arraybackslash}p{(\linewidth - 4\tabcolsep) * \real{0.3333}}
  >{\raggedright\arraybackslash}p{(\linewidth - 4\tabcolsep) * \real{0.3333}}
  >{\raggedright\arraybackslash}p{(\linewidth - 4\tabcolsep) * \real{0.3333}}@{}}
\toprule\noalign{}
\begin{minipage}[b]{\linewidth}\raggedright
Métrica
\end{minipage} & \begin{minipage}[b]{\linewidth}\raggedright
Expressão
\end{minipage} & \begin{minipage}[b]{\linewidth}\raggedright
Pergunta
\end{minipage} \\
\midrule\noalign{}
\endhead
Sensibilidade ou recall & \(VP/(VP+FN)\) & Quantos positivos
encontramos? \\
Especificidade & \(VN/(VN+FP)\) & Quantos negativos reconhecemos? \\
Precisão & \(VP/(VP+FP)\) & Quantos alertas estão corretos? \\
\bottomrule\noalign{}
\end{longtable}
}

Reduzir o cutoff aumenta ou mantém a sensibilidade e reduz ou mantém a
especificidade, no mesmo conjunto de escores. Precisão depende também da
prevalência.
\end{frame}

\begin{frame}{Desbalanceamento: acurácia pode esconder falhas}
\protect\phantomsection\label{desbalanceamento-acuruxe1cia-pode-esconder-falhas}
Se apenas 2\% das observações pertencem à classe 1, prever sempre 0
produz \textbf{98\% de acurácia}, mas \textbf{sensibilidade zero}.

Duas outras medidas são

\[
\operatorname{Balanced\ Accuracy}
=\frac{\operatorname{Sensibilidade}+\operatorname{Especificidade}}{2},
\]

\[
F_1=\frac{2VP}{2VP+FP+FN}.
\]

A primeira dá peso igual ao desempenho nas duas classes. A segunda
combina precisão e recall, sem usar diretamente os verdadeiros
negativos.

\textbf{A métrica deve refletir o objetivo. Nenhuma dessas medidas
representa automaticamente o custo real da decisão.}
\end{frame}

\begin{frame}{ROC, AUC e Precision--Recall}
\protect\phantomsection\label{roc-auc-e-precisionrecall}
Ao variar o cutoff, obtemos:

{\def\LTcaptype{none} % do not increment counter
\begin{longtable}[]{@{}lll@{}}
\toprule\noalign{}
Curva & Eixo horizontal & Eixo vertical \\
\midrule\noalign{}
\endhead
ROC & \(FPR=1-\text{especificidade}\) & \(TPR=\text{sensibilidade}\) \\
Precision--Recall & Recall & Precisão \\
\bottomrule\noalign{}
\end{longtable}
}

A AUC-ROC mede ordenação:

\[
\operatorname{AUC}=P(S_1>S_0)+\tfrac12P(S_1=S_0),
\]

com escores de observações positivas e negativas independentes. A
parcela final trata empates.

\textbf{AUC não avalia calibração nem determina cutoff.} Em classes
raras, a curva PR evidencia a proporção de alertas corretos e depende da
prevalência.
\end{frame}

\begin{frame}{As duas curvas do mesmo classificador}
\protect\phantomsection\label{as-duas-curvas-do-mesmo-classificador}
\pandocbounded{\includegraphics[keepaspectratio]{aula_altas_dimensoes_est335_revisada_V2_files/figure-beamer/unnamed-chunk-2-1.png}}

Escores simulados com 35\% de positivos. A linha horizontal em PR marca
essa prevalência. As curvas descrevem ordenação, sem garantir que os
escores sejam probabilidades calibradas.
\end{frame}

\begin{frame}{Três intervenções que precisam de validação}
\protect\phantomsection\label{truxeas-intervenuxe7uxf5es-que-precisam-de-validauxe7uxe3o}
{\def\LTcaptype{none} % do not increment counter
\begin{longtable}[]{@{}
  >{\raggedright\arraybackslash}p{(\linewidth - 4\tabcolsep) * \real{0.3333}}
  >{\raggedright\arraybackslash}p{(\linewidth - 4\tabcolsep) * \real{0.3333}}
  >{\raggedright\arraybackslash}p{(\linewidth - 4\tabcolsep) * \real{0.3333}}@{}}
\toprule\noalign{}
\begin{minipage}[b]{\linewidth}\raggedright
Intervenção
\end{minipage} & \begin{minipage}[b]{\linewidth}\raggedright
O que modifica
\end{minipage} & \begin{minipage}[b]{\linewidth}\raggedright
Exemplos
\end{minipage} \\
\midrule\noalign{}
\endhead
Dados de treinamento & Composição da amostra & Oversampling,
undersampling, SMOTE \\
Função objetivo & Peso dos erros no ajuste & Pesos de classes e
custos \\
Decisão final & Conversão do escore em ação & Alteração do cutoff \\
\bottomrule\noalign{}
\end{longtable}
}

Reamostrar não equivale a mudar apenas o cutoff. Pode alterar o modelo e
a interpretação de suas probabilidades.

\textbf{Dentro de cada divisão de validação:} ajustar seleção,
padronização e eventual reamostragem apenas no treino. Escolher
hiperparâmetros e cutoff usando treino/validação. Preservar o teste e
sua distribuição para a avaliação final.
\end{frame}

\begin{frame}{Combinação de classificadores}
\protect\phantomsection\label{combinauxe7uxe3o-de-classificadores}
\textbf{Hard voting:} combinar as classes previstas por maioria, com
regra de desempate definida.

\[
g_{\mathrm{ens}}(\mathbf x)=\operatorname{moda}\{g_1(\mathbf x),\ldots,g_M(\mathbf x)\}.
\]

\textbf{Soft voting:} combinar probabilidades de classes na mesma
escala:

\[
\bar p_k(\mathbf x)=\sum_{m=1}^{M}w_m\hat p_{mk}(\mathbf x),
\quad w_m\geq0,\quad \sum_mw_m=1.
\]

A combinação pode ajudar quando os erros dos modelos diferem. Pesos e
escolha de componentes pertencem à validação. Fazer a média não garante
calibração nem melhor desempenho.

Com custos iguais, escolher a maior probabilidade é natural. Com custos
assimétricos, aplicar a regra de decisão correspondente.
\end{frame}

\begin{frame}{Interpretar em alta dimensão}
\protect\phantomsection\label{interpretar-em-alta-dimensuxe3o}
\textbf{Interpretação global:} descrever o comportamento geral do
modelo.

\textbf{Interpretação local:} compreender uma predição específica.

Em alta dimensão, três cuidados são centrais:

\begin{itemize}
\tightlist
\item
  Entradas redundantes podem substituir umas às outras no ajuste.
\item
  Um modelo esparso é compacto, mas a seleção pode variar entre
  amostras.
\item
  Explicar a predição de um modelo não identifica um efeito causal.
\end{itemize}

Nos próximos slides, \(\hat g\) representa a saída numérica que queremos
explicar: resposta prevista, escore ou probabilidade de uma classe
especificada.
\end{frame}

\begin{frame}{PDP: comportamento médio do modelo}
\protect\phantomsection\label{pdp-comportamento-muxe9dio-do-modelo}
Para estudar a entrada \(X_j\), substituímos seu valor por \(t\) em
todas as observações e calculamos a média das predições:

\[
\boxed{\hat h_j(t)=\frac1N\sum_{i=1}^{N}\hat g(t,\mathbf x_{i,-j}).}
\]

\(\mathbf x_{i,-j}\) contém as demais entradas da observação \(i\).
\(N\) é o tamanho da base usada como referência para a explicação.

\textbf{Leitura:} se o PDP cresce, as predições médias do modelo
aumentam nas combinações consultadas.

Isso não é, em geral, \(\mathbb E(Y\mid X_j=t)\) nem um efeito causal.
Com entradas dependentes, pode consultar combinações pouco representadas
nos dados.
\end{frame}

\begin{frame}{ICE: comportamento por observação}
\protect\phantomsection\label{ice-comportamento-por-observauxe7uxe3o}
Em vez de calcular primeiro a média, traçamos uma curva para cada
observação:

\[
\hat h_{j,i}(t)=\hat g(t,\mathbf x_{i,-j}).
\]

A média dessas curvas, na mesma base de referência, é o PDP:

\[
\boxed{\hat h_j(t)=\frac1N\sum_{i=1}^{N}\hat h_{j,i}(t).}
\]

Curvas com inclinações ou formas diferentes mostram heterogeneidade no
comportamento do modelo e podem revelar interações.

\textbf{ICE também pode criar combinações pouco plausíveis quando
alteramos uma entrada e mantemos as correlacionadas fixas.}
\end{frame}

\begin{frame}{A média pode esconder heterogeneidade}
\protect\phantomsection\label{a-muxe9dia-pode-esconder-heterogeneidade}
\pandocbounded{\includegraphics[keepaspectratio]{aula_altas_dimensoes_est335_revisada_V2_files/figure-beamer/unnamed-chunk-3-1.png}}

Modelo ilustrativo \(\hat g(x_1,x_2)=x_1x_2\), sem ajuste a dados. O PDP
horizontal esconde a interação: o sentido da variação depende de
\(x_2\).
\end{frame}

\begin{frame}{LIME: aproximação local da predição}
\protect\phantomsection\label{lime-aproximauxe7uxe3o-local-da-prediuxe7uxe3o}
Para explicar a saída em \(\mathbf x^*\), o LIME:

\begin{enumerate}
\tightlist
\item
  Gera perturbações na vizinhança da observação.
\item
  Consulta o modelo original nessas entradas.
\item
  Ajusta um modelo simples, dando maior peso às perturbações próximas.
\end{enumerate}

\[
q^*=\arg\min_{q\in\mathcal G}
\{\mathcal L(\hat g,q,\pi_{\mathbf x^*})+\Omega(q)\}.
\]

\(\mathcal L\) mede fidelidade local, \(\pi_{\mathbf x^*}\) define os
pesos e \(\Omega\) penaliza a complexidade da explicação.

\textbf{A explicação descreve uma aproximação do modelo perto de um
ponto.} Ela depende da vizinhança, da geração de perturbações e da
representação utilizada.
\end{frame}

\begin{frame}{Redundância e explicações pouco confiáveis}
\protect\phantomsection\label{redunduxe2ncia-e-explicauxe7uxf5es-pouco-confiuxe1veis}
Imagine entradas quase equivalentes, como duas medições do mesmo sensor.

\textbf{PDP e ICE:} variar uma mantendo a outra fixa pode gerar
situações fora da região observada.

\textbf{LIME:} perturbar as duas independentemente pode criar vizinhos
artificiais. Muitos preditores também tornam a noção de proximidade
menos informativa.

\textbf{Lasso:} selecionar uma medição e descartar outra não prova que
apenas a selecionada tenha importância substantiva.

Antes de interpretar, examine suporte dos dados, estabilidade e
fidelidade. Uma explicação convincente não substitui desempenho fora da
amostra.
\end{frame}

\begin{frame}{Síntese integrada}
\protect\phantomsection\label{suxedntese-integrada}
{\def\LTcaptype{none} % do not increment counter
\begin{longtable}[]{@{}
  >{\raggedright\arraybackslash}p{(\linewidth - 2\tabcolsep) * \real{0.5000}}
  >{\raggedright\arraybackslash}p{(\linewidth - 2\tabcolsep) * \real{0.5000}}@{}}
\toprule\noalign{}
\begin{minipage}[b]{\linewidth}\raggedright
Etapa
\end{minipage} & \begin{minipage}[b]{\linewidth}\raggedright
Pergunta central
\end{minipage} \\
\midrule\noalign{}
\endhead
Estimar em alta dimensão & Que estrutura reduz a complexidade
relevante? \\
Decidir & Que ação minimiza o custo esperado? \\
Avaliar & Que métricas representam o objetivo em dados novos? \\
Interpretar & A explicação é fiel e usa combinações plausíveis? \\
\bottomrule\noalign{}
\end{longtable}
}

\textbf{Esparsidade e redundância afetam todo o processo:} a estimação,
a estabilidade das variáveis selecionadas e a confiabilidade das
explicações.

O objetivo é obter predições úteis, decisões justificadas e explicações
compatíveis com o que os dados permitem concluir.
\end{frame}

\begin{frame}{Referências}
\protect\phantomsection\label{referuxeancias}
\begin{itemize}
\tightlist
\item
  \href{https://www.professortp.com.br/est335/aulas.html}{Aulas da
  EST335}: regularização, métodos locais, SVM e redes neurais.
\item
  Material de apoio: \emph{Decisão, Avaliação e Interpretabilidade},
  EST335, arquivo fornecido para esta aula.
\item
  James et al.~(2021). \emph{An Introduction to Statistical Learning},
  2ª ed.~\href{https://www.statlearning.com/}{Livro}.
\item
  Wasserman. \emph{Nonparametric Regression}. Notas de Statistical
  Machine Learning.
  \href{https://stat.cmu.edu/~larry/=sml/nonpar.pdf}{Notas}.
\item
  Gautier e Tsybakov (2016). \emph{Regularity Properties for Sparse
  Regression}.
  \href{https://pmc.ncbi.nlm.nih.gov/articles/mid/NIHMS750383/}{Artigo}.
\end{itemize}
\end{frame}




\end{document}
